\section{A Recapitulation of the Dynamical Vocabulary}
\label{sec:dynsys}

\noindent
The paper describes the lifecycle of a shared flourishing in the vocabulary of dynamical systems, and the central terms are recalled here in a form sufficient for the use made of them. The recapitulation is brief and non-technical, and it carries forward the qualification, Standard treatments of the notions recalled here are \citet{strogatz1994} and \citet{hirsch2013}. established at length in the prior paper, that the relation is modelled by a dynamical system only as a structural analogy, and that the analogy is subject to the limit there identified, that the law governing a relation is itself reconfigured by what the relation generates, so that no fixed system captures it in full. The vocabulary is employed for what it makes visible, and not as a claim that a relation is a system of the standard kind.

\subsection{States, evolution, and trajectories}
\label{sec:dyn-states}

A dynamical system is specified by a space of states, the configurations the system may occupy, together with a rule of evolution that carries each state to its successor in time. The history of the system is a trajectory through the state space, the path traced as the state evolves under the rule. For a relation, the state space is the space of configurations a shared life may occupy, and the trajectory is the relation's history; the rule of evolution is the relation's own dynamics, which, as the prior paper insisted, is not fixed but is reconfigured by the relation's products, a qualification held in reserve throughout what follows.

\subsection{Fixed points, stability, and decay}
\label{sec:dyn-fixed}

A fixed point is a state that the rule of evolution carries to itself, a configuration at which the system, once arrived, remains. A fixed point is stable if trajectories that begin nearby are drawn toward it, and unstable if they are driven away. A trajectory that loses the energy or the input sustaining its movement winds down toward a fixed point and there comes to rest. For a relation, the pertinent fixed point is the state in which the relation's movement has ceased, in which value is no longer generated, exchanged, or returned, and the parties no longer revise their sense of one another. This terminal rest is not a benign equilibrium but the cessation of the relation as a living process, and the paper names it extinction (\xref{sec:regimes}). A relation that is not sustained decays toward it.

\subsection{Attractors and limit cycles}
\label{sec:dyn-attractors}

An attractor is a set of states toward which trajectories from a region of the state space are drawn, and on which they thereafter remain. A fixed point is the simplest attractor. A limit cycle is an attractor of a different kind: a closed trajectory, a loop in the state space, toward which neighbouring trajectories are drawn and along which the system then circulates indefinitely. A limit cycle is a self-sustaining oscillation, a movement that maintains itself by returning, again and again, through a cycle of states rather than coming to rest at a point. The relevance to the present paper is direct: a living relation is modelled not by a fixed point, which would be its cessation, but by a sustained cyclic movement, the lifecycle of cognition, practice, evaluation, reflection, and renewed cognition, which maintains the relation by its continued turning. The persistence of this cyclic movement, rather than the attainment of any resting state, is the mark of a living relation.

\subsection{The phase accumulated around a cycle}
\label{sec:dyn-phase}

A system that returns, around a cycle, to its starting configuration in the state space need not return unchanged in every respect. A quantity carried along the cycle may, on return to the starting configuration, fail to coincide with its initial value, and the discrepancy, the phase accumulated around the loop, is the holonomy on which the prior papers turned. A cycle of vanishing holonomy returns everything to where it began, a pure repetition that goes round without rising. A cycle of positive holonomy returns to the starting configuration with a gain, a spiral that ascends as it returns. A cycle of negative holonomy returns with a loss, a winding that depletes as it turns. The distinction between the good and the bad cycle, developed throughout the series, is the sign of this accumulated phase, and it is carried over here to the lifecycle of a shared flourishing: the same five moments may be traversed in a manner that ascends, in a manner that merely repeats, or in a manner that depletes, and the difference is the holonomy of the turn, not the identity of the moments traversed.

\subsection{Bifurcation}
\label{sec:dyn-bifurcation}

The qualitative behaviour of a dynamical system may depend on parameters, the slowly varying conditions under which the system evolves. As a parameter is varied, the system's behaviour may change smoothly through a range and then, at a critical value, change in kind: a stable fixed point may lose its stability, a limit cycle may appear or vanish, an attractor may divide. Such a qualitative change at a critical parameter value is a bifurcation. The relevance is that the regime of a relation, whether it sustains a flourishing spiral, settles into a depleting cycle, or decays toward extinction, may not change gradually with its conditions but may shift in kind at a threshold, a small further change in the sustaining conditions carrying the relation across a critical value into a qualitatively different regime. The language of bifurcation is the language in which the paper describes the passage of a relation from one regime to another (\xref{sec:regimes}), and it captures a feature the gradual language of more or less well-being does not, that the difference between a sustained and a collapsing relation can be a difference in kind reached by a small change in degree.
